\subsubsection{Influencia de la estructura del contexto en el razonamiento cronológico}

En la notebook 9 se experimenta con la forma en que la estructura del contexto afecta el razonamiento cronológico y la extracción de datos por parte de modelos de lenguaje como GPT-3.5. A continuación, se analizan los pasos principales del experimento realizado en \href{https://github.com/aditzend/hyperpersonalization/blob/main/app/notebooks/9.ipynb}{9.ipynb}.

\paragraph{Configuración del entorno experimental}

Se utiliza FAISS para la búsqueda semántica y OpenAI Embeddings para vectorizar los mensajes. Se configura también una función que interactúa con los modelos de lenguaje a través de la API de OpenAI.

La transcripción íntegra de este listado figura en el \textit{Anexo técnico} (véase el Apéndice~\ref{anexo:code-4-1-9-1}).

\paragraph{Contextos condensados y simulación de conversaciones}

Se simulan varias conversaciones entre el usuario y el sistema. En una de ellas, el usuario pide un turno para su madre, llamada Alba, y menciona que tiene 66 años. En otra, el usuario pregunta la capital de Colombia.

\paragraph{Ejemplos de mensajes de la conversación principal:}
\begin{itemize}
    \item Usuario: ``Quiero un turno para mi madre''.
    \item Asistente: ``¿Cuál es el nombre de tu madre?''.
    \item Usuario: ``Alba''.
    \item Asistente: ``¿Qué edad tiene?''.
    \item Usuario: ``66''.
    \item Asistente: ``El turno es para el 30 de agosto a las 15:00''.
\end{itemize}

La transcripción íntegra de este listado figura en el \textit{Anexo técnico} (véase el Apéndice~\ref{anexo:code-4-1-9-2}).

\paragraph{Extracción de datos con GPT-3.5}

En lugar de simplemente almacenar las conversaciones, se utiliza el modelo GPT-3.5 para extraer la información clave de cada conversación. Esto incluye detalles como el nombre y la edad de la madre del usuario, así como la fecha y hora en que la información fue recopilada. Esta extracción se almacena en FAISS junto con los metadatos de la conversación.

La transcripción íntegra de este listado figura en el \textit{Anexo técnico} (véase el Apéndice~\ref{anexo:code-4-1-9-3}).

\paragraph{Ejemplo de extracción obtenida:}
``El usuario mencionó que quiere un turno para su madre. La madre del usuario se llama Alba y tiene 66 años.''

\paragraph{Consultas simples y razonamiento cronológico}

Se comienza con una consulta simple: ``¿Qué edad tiene Alba?''. El sistema encuentra el contexto relevante y responde correctamente: ``Alba tiene 66 años''.

La transcripción íntegra de este listado figura en el \textit{Anexo técnico} (véase el Apéndice~\ref{anexo:code-4-1-9-4}).

Luego, se plantea una consulta más compleja: ``¿Qué edad tenía Alba en 2020?''. En este caso, GPT-3.5 logra razonar correctamente y responde que Alba tenía 63 años en 2020, basándose en la edad actual de 66 años.

\paragraph{Consultas adicionales sobre fechas}

Se continúa con una pregunta más difícil: ``¿Qué edad tenía Alba en 1970?''. El sistema nuevamente razona correctamente y responde que Alba tenía 13 años en 1970, lo que demuestra que, con la forma correcta de estructurar y almacenar el contexto, GPT-3.5 es capaz de realizar razonamientos cronológicos precisos.

La transcripción íntegra de este listado figura en el \textit{Anexo técnico} (véase el Apéndice~\ref{anexo:code-4-1-9-5}).

\textbf{Resultados obtenidos:}
\begin{itemize}
    \item ``¿Qué edad tiene Alba?'' → ``Alba tiene 66 años.''
    \item ``¿Qué edad tenía Alba en 2020?'' → ``Alba tenía 63 años en 2020.''
    \item ``¿Qué edad tenía Alba en 1970?'' → ``La edad de Alba en 1970 sería 13 años.''
\end{itemize}

\paragraph{Consulta sobre reconocimiento de datos personales}

Finalmente, se realiza una consulta general: ``¿Conoces a mi madre?''. El sistema responde de manera acertada, reconociendo a Alba y proporcionando los detalles previamente almacenados: ``Sí, conozco a tu madre. Su nombre es Alba y tiene 66 años''.

La transcripción íntegra de este listado figura en el \textit{Anexo técnico} (véase el Apéndice~\ref{anexo:code-4-1-9-6}).

\paragraph{Conclusiones de la novena iteración}

La notebook 9 demuestra que la estructura del contexto influye significativamente en la capacidad del modelo para realizar razonamientos cronológicos y responder a consultas sobre datos específicos. Al extraer y almacenar los datos de forma clara y precisa, el modelo es capaz de utilizar esta información para responder correctamente incluso a preguntas complejas sobre fechas y edades, mejorando su rendimiento en tareas de razonamiento temporal.

Los resultados evidencian que:
\begin{itemize}
    \item La extracción de información estructurada mejora la precisión de las respuestas
    \item GPT-3.5 puede realizar cálculos cronológicos cuando el contexto está bien organizado
    \item La vectorización semántica permite recuperar eficientemente información relevante
    \item El sistema mantiene coherencia en el reconocimiento de entidades personales a lo largo del tiempo
\end{itemize} 